% 第 5 章：RA 式分阶段灾害 MILP
\section{RA 式分阶段灾害 MILP}\label{sec:res-stage}

\subsection{功能定位}
\fld{RAStyleStageMILP}（\srcpath{run\_distribution\_resilience\_stage\_milp\_assessment}，
\srcpath{include/hacdcpf/resilience/resilience\_assessment.hpp:run\_distribution\_resilience\_stage\_milp\_assessment}；由 \fld{use\_ra\_style\_stage\_milp} 或首选开关
\fld{enable\_disaster\_stages} 触发）以“灾害隔离 + 灾后重构”两阶段、开关约束的拓扑 MILP
求解，阶段语义见 \fld{DistributionDisasterStage}（:DistributionDisasterStage）：
\texttt{Normal} $\to$ \texttt{DisasterIsolation} $\to$ \texttt{DisasterPostFaultReconfig}
$\to$ \texttt{PostDisasterRepair}。

\subsection{分阶段模型}
\begin{itemize}
  \item \textbf{阶段 1 灾害隔离}：以开关隔离受灾支路
        （\fld{use\_switch\_based\_fault\_isolation}=true，\fld{allow\_stage1\_open\_switches}），
        非故障支路的操作须有开关（\fld{require\_switch\_for\_nonfault\_branch\_operation}），
        默认禁止无开关直接操作支路（\fld{allow\_branch\_operation\_without\_switch}=false）；
  \item \textbf{阶段 2 灾后重构}：在 \fld{post\_fault\_reconfig\_window\_hr}（默认 2h）内闭合联络
        （\fld{allow\_stage2\_close\_ties}）以恢复负荷，可限 \fld{use\_remote\_switch\_only}；
  \item 每阶段使用有功输运、节点供需、带电与径向拓扑约束；不建立严格多时段 MIP 的 AC
        支路额定流限或电压包络，故相应 validity flags 为 false。
\end{itemize}
MESS 既可由专用严格 MILP（\fld{use\_strict\_mip\_for\_mess}）也可阶段内调度并在失败时回退
（\fld{fallback\_to\_stage\_mess\_dispatch}）。逐步结果记灾害阶段 \fld{disaster\_stage}、隔离性开断
\fld{isolation\_open\_ac\_branch\_ids}/\fld{isolation\_open\_dc\_branch\_ids} 与重构性合断，并区分
\fld{isolation\_switch\_actions} 与 \fld{reconfiguration\_switch\_actions}。

\subsection{与其他模型的关系}
三种模型互斥选择、共享同一逐步结果结构与弹性指标；启发式最快、多时段 MILP 全局协调、
RA 式分阶段贴合灾害隔离-重构的工程流程。三者的能力边界均由 \fld{ValidityFlags} 显式声明，
拓扑求解失败显式反映；MESS 子问题若允许回退则在 \fld{mess\_dispatch\_model} 与状态中明确标识。
